\section{Albanese fibre groups and completion of the proof}\label{sec:completion}

The remaining task is group-theoretic and geometric assembly: isolate
the Albanese fibre-group image, establish the minimality hypotheses
of the positivity theorem, and apply torus descent in both
representation cases.

Let \(X\) be special. If \(\dim X=0\), its fundamental group is trivial,
so assume \(\dim X>0\). Every connected finite \'etale cover is special
and has surjective Albanese morphism by Theorem~\ref{thm:specialness-inputs}.
Its first Betti number is therefore at most \(2\dim X\). Pass to a cover
on which this number is maximal, and continue to denote the cover by
\(X\). Pullback on \(H^1(-,\mathbb Q)\) is injective by transfer, so
any further connected finite \'etale cover has the same first Betti
number. Write
\[
 \alpha:X\longrightarrow A=\Alb(X),\qquad
 G=\pi_1(X),\qquad N=\im\bigl(\pi_1(S)\longrightarrow G\bigr),
\]
where \(S\) is a smooth general Albanese fibre. The group \(N\) is
finitely generated and normal. When \(A\) is a point, these statements
mean \(S=X\) and \(N=G\).

\subsection{Exactness over an arbitrary Albanese torus}

\begin{lemma}\label{lem:albanese-exactness}
The natural sequence is exact:
\[
 1\longrightarrow N\longrightarrow G
 \xrightarrow{\alpha_*}\pi_1(A)\longrightarrow1.
\]
\end{lemma}
\begin{proof}
Set \(R=G/N\). Over the smooth fibration locus in \(A\), fibrewise
exactness gives a monodromy homomorphism to \(R\), compatible with the
quotient map on the total space. A meridian around a base divisor has
finite order in \(R\), since a component of its inverse image kills a
positive power of that meridian. Form the actual inertial divisor
\(D_A\) using these orders, as in Section~\ref{sec:descent}.

Suppose \(D_A\ne0\). Lemma~\ref{lem:torus-divisor} gives a quotient
torus \(p:A\to B\) with connected fibres, where \(B\) is an abelian
variety of dimension \(b>0\), and an ample effective rational divisor
\(D_B\) such that \(D_A=p^*D_B\). The smooth quotient map pulls back
each prime divisor with multiplicity one, so the components of \(D_B\)
retain the meridian orders of their inverse images.

Lemma~\ref{de:inertia-divisibility}, applied to the smooth torus
bundle \(A\to B\), shows that each such order divides the
multiplicity of \emph{every} prime divisor of \(X\) dominating the
corresponding divisor of \(B\) under \(p\alpha\). This includes
components whose images in \(A\) have larger codimension. For prime
divisors whose images in \(B\) have codimension at least two, the
top differential of \(p\alpha\) vanishes: at their general points
its restriction to the divisor has rank at most \(b-2\), so the
total rank is at most \(b-1\). There are only finitely many relevant
components over the fixed support of \(D_B\). Thus a sufficiently
small rational \(\varepsilon>0\) makes their differential vanishing
orders compensate the poles allowed by \(\varepsilon D_B\).
This is the fixed-map argument of Lemma~\ref{de:small-boundary}.

Let \(L\subset\Omega_X^b\) be the saturated line generated by
the top differentials from \(B\). For sufficiently divisible \(m\),
pullback of pluriforms gives
\[
 H^0\bigl(B,m(K_B+\varepsilon D_B)\bigr)
 \longrightarrow H^0(X,L^{\otimes m}).
\]
The pullback is injective, and ratios of sections retain their
independence because \(X\to B\) is surjective. Since \(K_B\) is
trivial and \(\varepsilon D_B\) is ample, \(\kappa(X,L)\ge b\).
The opposite inequality follows from
Theorem~\ref{thm:specialness-inputs}. This contradicts specialness.

All meridians therefore have trivial image in \(R\). The monodromy
homomorphism factors through \(\pi_1(A)\), by normal meridian
generation for the complement of the discriminant. Its composition
with \(G\to\pi_1(A)\) is the quotient \(G\to R\). Conversely,
\(\alpha_*\) kills \(N\) and factors through \(R\).
These two factorizations are inverse because the maps from \(G\)
are surjective. This proves \(R\cong\pi_1(A)\).
\end{proof}

For a further connected finite cover \(X'\to X\) corresponding to
\(G'\le G\), write \(A'=\Alb(X')\). The induced homomorphism
\(A'\to A\) is surjective and, by maximality of the first Betti
number, is an isogeny. Lemma~\ref{lem:albanese-exactness} on both
manifolds shows that the new Albanese fibre-group image is
\begin{equation}\label{eq:fibre-intersection}
 N'=N\cap G'.
\end{equation}
Indeed the map on torus fundamental groups induced by an isogeny
is injective. Lemma~\ref{lem:realize-finite-index} consequently
realizes every prescribed finite-index condition on \(N\) after
shrinking the total group by finite index.

\subsection{Finite abelianization on all finite covers}

\begin{proposition}\label{prop:fibre-fab}
The group \(N\) is FAb.
\end{proposition}
\begin{proof}
We first show that every group \(N'=N\cap G'\) arising from a
finite cover has \(\Hom(N',\C)=0\). Set
\(\Lambda=\pi_1(A')\) and
\(V=\Hom(N',\C)\). The vector space \(V\) is finite dimensional
because \(N'\) is finitely generated. Conjugation defines an action
of the abelian group \(\Lambda\) on it. Write \(T_\lambda\) for the
operator
\[
 (T_\lambda v)(n)=v(\widetilde\lambda^{-1}n\widetilde\lambda),
\]
where \(\widetilde\lambda\) is any lift to \(G'\). Inner conjugation
acts trivially on additive homomorphisms, so the operator is
independent of this lift.

For a nontrivial character \(\beta:\Lambda\to\C^*\), the
cohomology of \(\Lambda\) with coefficients in \(\C_\beta\)
vanishes in every degree. This follows, for example, from the
Koszul resolution: one generator acts by a scalar different from
one, making that factor contractible. Inflation--restriction gives
\begin{equation}\label{eq:twisted-inflation}
 H^1(G',\C_\beta)
 \cong (V\otimes\C_\beta)^\Lambda.
\end{equation}
Choose a basis \(\lambda_1,\ldots,\lambda_r\) of \(\Lambda\) and put
\(T_i=T_{\lambda_i}\). A nonzero invariant on the right of
\eqref{eq:twisted-inflation} corresponds to a vector \(w\ne0\)
satisfying
\[
 T_iw=\beta(\lambda_i)^{-1}w\qquad(1\le i\le r).
\]
Each \(T_i\) has finitely many eigenvalues, and the values on this
basis determine a character. Hence only finitely many nontrivial
\(\beta\) have \(H^1(G',\C_\beta)\ne0\).

Inflation identifies the character group of \(\Lambda\) with the
open and closed identity component of \(\operatorname{Hom}(G',\C^*)\):
the Albanese lattice is \(H_1(X',\mathbb Z)\) modulo torsion, and the
character group has finitely many components. Group and manifold
cohomology agree in degree one for these local systems. The nontrivial
support points just found are therefore isolated in the full
degree-one rank-one cohomology jump locus. B.~Wang's torsion-translate
theorem for compact K\"ahler manifolds
\cite[author's version, Theorem~1.3]{BWang2016} makes them torsion.

Suppose \(V\ne0\). The commuting complex operators \(T_i\) have a
common eigenvector \(v\ne0\): successively take an eigenspace of the
next operator inside the nonzero common eigenspace already chosen,
which remains stable by commutativity. Their invertibility makes the
eigenvalues nonzero, and they define a character
\(\chi:\Lambda\to\C^*\). If \(\chi\ne1\), the vector \(v\) makes
\(\beta=\chi^{-1}\) one of the nontrivial support points above, so
\(\chi\) is torsion. If \(\chi=1\), it is already torsion. For
\(r=0\), take any nonzero \(v\) and the unique trivial character.

Let \(A_0\to A'\) be the isogeny with lattice
\(\Lambda_0=\ker\chi\), and pull it back to \(X'\). The total-space
cover \(X_0\) is connected because \(G'\to\Lambda\) is surjective;
its group is \(G_0=(\alpha'_*)^{-1}(\Lambda_0)\), with kernel still
exactly \(N'\). Thus the vector space \(V=\Hom(N',\C)\) has not
changed, and the same \(v\) is fixed by the restricted lattice.
The pulled-back fibration has connected fibres. By the Albanese
universal property it factors through a surjective homomorphism
\(\Alb(X_0)\to A_0\), which maximality makes an isogeny.
Connectedness of the pulled-back fibres forces its degree to be one.
We may therefore relabel \(X_0,A_0,G_0,\Lambda_0\) as
\(X',A',G',\Lambda\), with the same \(V\) and \(V^\Lambda\ne0\).

The untwisted inflation--restriction sequence contains
\begin{equation}\label{eq:transgression}
 \begin{aligned}
 0&\longrightarrow H^1(\Lambda,\C)
 \longrightarrow H^1(G',\C)
 \longrightarrow V^\Lambda\\
 &\longrightarrow H^2(\Lambda,\C)
 \longrightarrow H^2(G',\C).
 \end{aligned}
\end{equation}
The last arrow is injective. To see this, its composition with
the natural map to \(H^2(X',\C)\) is the pullback
\(H^2(A',\C)\to H^2(X',\C)\), since a torus is a
\(K(\Lambda,1)\). That pullback is injective for a surjective
map from a compact K\"ahler manifold. In fact, if
\(a=\dim A'\) and \(d=\dim X'-a\), fibre integration of a
K\"ahler class \([\omega]^d\) is a positive constant. For
\(\eta\in H^2(A',\C)\) and \(\theta\in H^{2a-2}(A',\C)\),
\[
 \int_{X'}\alpha'^*(\eta\smile\theta)\smile[\omega]^d
 =c\int_{A'}\eta\smile\theta,\qquad c>0.
\]
Poincar\'e duality proves the claim. For \(a=0\) there is no
degree-two group to consider.

The transgression in \eqref{eq:transgression} is zero.
Hence a nonzero invariant vector would give
\(b_1(X')>2\dim A'\), contrary to the defining equality for
the Albanese torus. The isogeny kept \(N'\) and \(V\) unchanged, so
this contradiction proves \(V=0\) for the original cover as well.
Maximality ensured that all intervening Albanese tori were isogenous
and that the pulled-back torus was itself the new Albanese torus.

Finally, let \(N_0\le N\) have finite index. By
Lemma~\ref{lem:realize-finite-index}, it contains some
\(N\cap G'\). The latter has no nonzero homomorphism to \(\C\).
A homomorphism from \(N_0\) to \(\C\) restricting to zero on
a finite-index subgroup is zero. Thus \(\Hom(N_0,\C)=0\).
Since \(N_0\) is finitely generated, its abelianization is finite.
\end{proof}

If \(N\) is finite, Lemmas~\ref{lem:albanese-exactness}
and~\ref{lem:finite-by-abelian} prove the theorem. For the rest
of the proof suppose that \(N\) is infinite.

\subsection{A simultaneous choice of general Albanese fibres}

We specify the generality needed to globalize a fibration from a
single Albanese fibre. A finitely generated group has countably
many finite-index subgroups. Thus there are countably many connected
finite covers of \(X\), up to the choices relevant here. Their
Albanese tori are isogenous to \(A\).

On each cover, Barlet's space of compact analytic cycles
\cite{Barlet1975} has countably many irreducible components, each
compact in the K\"ahler case
\cite[Section~3, Convention~3.1]{CampanaGamma}.
We will use a countable list of families with holomorphic maps
\[
 p_j:W_j\longrightarrow T_j,\qquad e_j:W_j\longrightarrow X,\qquad
 W_j^\circ=p_j^{-1}(T_j^\circ)\subset W_j.
\]
Here \(W_j\) is a connected smooth compact complex manifold,
\(T_j^\circ\) is a connected smooth dense analytic Zariski open
subset of a compact parameter space \(T_j\), and
\(p_j:W_j^\circ\to T_j^\circ\) is proper and smooth with connected
fibres. Every irreducible reduced compact cycle \(V\subset X\)
occurs as \(e_j((W_j)_t)=V\) for some \(t\in T_j^\circ\), with the
fibre map proper and bimeromorphic. In particular,
\(W_j^\circ\) is the complement of a proper analytic subset in
the same compact \(W_j\).

Here is a construction of this list. On a fixed compact irreducible
component \(P\) of the space of \(k\)-cycles, K\"ahler volume is
constant, say \(v\). A cycle is reducible or nonreduced precisely
when it is a sum of two nonzero cycles. Both summands have volume at
most \(v\). Bounded-volume compactness and local finiteness of the
components of the Barlet space leave only finitely many possible
component pairs for these summands. Addition of cycles is holomorphic,
and its images from these compact pairs are analytic by proper mapping.
Thus the decomposable locus in \(P\) is a closed analytic subset.
It is proper whenever \(P\) contains an irreducible reduced cycle;
components containing no such cycle can be omitted.

Resolve \(P\) and restrict to the locus where this resolution is an
isomorphism. Pull back the reduced universal incidence, take the
closure of its component dominating the irreducible-cycle locus,
and resolve that compact incidence. Delete the critical parameter
values, the
parameters whose whole support fibre lies in the exceptional image,
and those for which a component of the full non-isomorphism locus
has full fibre dimension.
These are proper analytic subsets by proper generic smoothness and
the proper fibre-dimension theorem. Over the remaining open, each
resolved fibre is smooth, connected, and bimeromorphic to its cycle.
Repeat on the irreducible components of the omitted parameter subsets,
including the image omitted in resolving \(P\). Each repetition lowers
the parameter dimension, and a compact analytic subset has finitely
many irreducible components. The countably many initial components
therefore give the claimed countable list. This construction requires
no K\"ahler metric on \(W_j\) or \(T_j\); their role is properness
and fundamental-group transport.

Each nondominant evaluation image \(e_j(W_j)\) is a proper analytic
subset of the cover. On the smooth Albanese locus let \(s\) be the
fibre dimension. An evaluation image dominating the torus has
general fibre dimension at most \(s-1\); the locus where this
dimension is at least \(s\) is proper analytic by the proper
fibre-dimension theorem. An image not dominating the torus already
projects to a proper analytic subset. Exclude these loci and the
singular-fibre locus. A smooth Albanese fibre contained in an evaluation image
has been excluded. Make these exclusions for every family and finite
cover, and project them to \(A\) through the finite isogenies.
The resulting union is countable. Call a point outside it \emph{good},
and use the same term for any of its lifts to a later Albanese torus.

If the torus has dimension zero, its fibre is the whole cover,
which no proper evaluation image contains; no exclusion is needed
for these images. Thus this convention includes irregularity zero.
For any fibre above a good point, every nondominant evaluation
image meets that fibre in a proper analytic subset. A countable
union of such subsets cannot cover an open part of that fibre.
This will allow us to choose a fibration cycle in a family
dominating the total space.

\subsection{The minimal relative base}

Consider all finite covers just described and their good Albanese
fibres \(S\). For each connected meromorphic fibration
\(S\dashrightarrow T\), let \(C\) be the image in the corresponding
Albanese fibre group \(N\) of a resolved general fibration fibre.
Then \(C\triangleleft N\). Among these fibrations for which
\(N/C\) is infinite, minimize \(\dim T\), and denote the minimum
by \(\ell\). The identity fibration is admissible because \(N\)
is infinite; the fibration to a point is not. Hence
\[
 1\le\ell\le\dim S.
\]
All covers used below remain within this simultaneous minimum.

\begin{proposition}\label{prop:minimal-base}
After a further connected finite \'etale cover, there is a
fibration \(c:X\dashrightarrow Z\) over the Albanese torus
\(A\), with connected general fibres and relative base dimension
\(\ell\). If \(B_0\) is the actual image of a general
\(c\)-fibre group, then \(B_0\subseteq N\), \(B_0\triangleleft G\),
and \(Q=N/B_0\) is infinite.

For each prepared smooth model of \(Z\), a very general good point
of \(A\) gives an Albanese fibre \(S\) and a torus fibre \(Y\).
With the restricted actual inertial boundary \(\Delta_Y\) for
\(R=G/B_0\), there is a surjection
\[
 \pi_1^\orb(Y,\Delta_Y)\twoheadrightarrow Q
\]
compatible with \(S\dashrightarrow Y\) and the natural
\(\nu:\pi_1(S)\twoheadrightarrow Q\). For every finite-index
subgroup \(Q_0\le Q\), there is a connected finite \'etale cover
\(S'\to S\) with induced homomorphism \(\nu'\) whose image
\(Q'=\im\nu'\) lies in \(Q_0\).
This one cover satisfies
\[
 \operatorname{gdim}(S',\psi\circ\nu')\ge\ell
\]
for every surjection \(\psi:Q'\twoheadrightarrow P\) with \(P\) infinite.
\end{proposition}
\begin{proof}
Choose a cover, a good fibre \(S\), and a fibration attaining
\(\ell\). Write \(C\triangleleft N\) for its fibre-group image.
The general fibration fibres give compact irreducible cycles of
dimension \(\dim S-\ell\), after taking their images in \(X\).
Choose one through a point of the fibration's general locus in
\(S\) outside all the nondominant evaluation images prepared
above. Choose a family from the list, with \(p:W\to T\) and
\(e:W\to X\), containing this member. Its compact evaluation must dominate \(X\),
and hence is surjective.

The homotopy sequence of the smooth proper family
\(W^\circ=p^{-1}(T^\circ)\to T^\circ\) makes its fibre-group image
normal in \(\pi_1(W^\circ)\). The selected fibre is bimeromorphic
to the selected cycle, so its image under \(e_*\) is \(C\).
Consequently \(H=e_*\pi_1(W^\circ)\) normalizes \(C\).
The complement of \(W^\circ\) in the smooth \(W\) is analytic,
so \(\pi_1(W^\circ)\to\pi_1(W)\) is surjective and
\(H=e_*\pi_1(W)\).

Put \(\Lambda=\pi_1(A)\) and
\(\Lambda_0=(\alpha e)_*\pi_1(W)\). Lift the surjective holomorphic
map \(\alpha e:W\to A\) to the connected complex covering
\(A_{\Lambda_0}\to A\). Its lifted image is compact analytic by
proper mapping. Surjectivity downstairs forces this image to have
dimension \(\dim A\), so it is all of the connected manifold
\(A_{\Lambda_0}\). The covering is therefore compact and has finite
degree. Thus \(\alpha_*(H)=\Lambda_0\) has finite index in \(\Lambda\).

Let \(L=N_G(C)\). It contains \(N\) and \(H\), so Albanese exactness
gives
\[
 [G:L]=[\Lambda:\alpha_*(L)]
 \le[\Lambda:\alpha_*(H)]<\infty.
\]
Pass to the cover corresponding to \(L\). Its fibre-group image is
still \(N\), by \eqref{eq:fibre-intersection}, and \(C\) is now normal
in the total group. Continue to write \(X,G,A,\Lambda\) for this
cover, its group, its Albanese torus and its lattice. The whole
compact evaluation \(e:W\to X\) lifts,
since \(e_*\pi_1(W)=H\subseteq L\). Its lift has full rank somewhere,
and its compact analytic image is therefore the entire connected cover.
The proper fibre-dimension theorem, applied to
\(W\setminus W^\circ\subsetneq W\), shows that a general evaluation
fibre meets \(W^\circ\). Hence members with parameters in \(T^\circ\)
cover general total points. Smooth proper transport gives their fibre-group image
\(C\) after compatible base-point transport.

Apply Theorem~\ref{thm:gamma-reduction} to \(G\to G/C\), giving
\(c:X\dashrightarrow Z\). The covering cycles just constructed
have dimension \(\dim S-\ell\) and are contracted through very
general points. Thus
\[
 \dim Z\le\dim X-(\dim S-\ell)=\dim A+\ell.
\]
This map is over \(A\): a compact
subvariety with finite image in the quotient lattice
\(G/N\) maps constantly to \(A\). To verify the last assertion,
resolve the subvariety. A finite image in the torsion-free lattice
is trivial, so its map lifts to the universal cover of \(A\),
a complex vector space; compactness makes that lift constant.
Consequently the Albanese map factors through \(c\), after
resolving. The induced map \(g:Z\to A\) has connected general
fibres: a nontrivial finite Stein factor of \(g\) would split a
general fibre of the connected Albanese map.

Let \(B_0\) be the actual group image of a general \(c\)-fibre.
It is normal in \(G\), lies in \(N\), and has finite image
\(B_0C/C\) in \(G/C\). If \(N/B_0\) were finite, the image of
\(N\) in \(G/C\) would be a finite union of translates of this
finite image, contradicting \(N/C\) infinite. Thus \(N/B_0\)
is infinite. For the dimension comparison, restrict \(c\) to a
very general good Albanese fibre. On simultaneous smooth loci,
its fibres are the same general \(c\)-fibres, so they are connected
and have image \(B_0\) in \(N\). This restricted fibration is
admissible in the minimum, giving \(\dim Z-\dim A\ge\ell\).
Equality follows. No equality or inclusion between \(B_0\) and
\(C\) has been used.

For inertia use \(R=G/B_0\), so
\[
 1\longrightarrow Q=N/B_0\longrightarrow R=G/B_0
 \longrightarrow\Lambda\longrightarrow1.
\]
Fix any prepared smooth model of \(Z\), with actual inertial
boundary \(D_Z\) for \(R\), and write \(\mu:\widetilde X\to X\)
for its resolved source, with \(c:\widetilde X\to Z\) holomorphic.
Choose afresh a very general good point
\(a\in A\) so that
\[
 S=\alpha^{-1}(a),\qquad
 \widetilde S=(\alpha\mu)^{-1}(a),\qquad Y=g^{-1}(a)
\]
are smooth, \(\widetilde S\to S\) is a modification, and \(Y\)
is transverse to every dominating stratum of the deleted divisor
and avoids the images of the other strata. This choice may depend
on the prepared model. The induced \(S\dashrightarrow Y\) is
meromorphic, and \(\dim Y=\ell\).

Let \(Z^\circ\) be the prepared smooth fibration locus for \(c\), put
\(Y^\circ=Y\cap Z^\circ\), and put
\(\widetilde S^\circ=c^{-1}(Y^\circ)\). The monodromy
\(\psi_Z:\pi_1(Z^\circ)\twoheadrightarrow R\) has composite
\(g_*\) in \(\Lambda\), by compatibility after precomposing with
the surjection from the open source group. Its restriction to
\(\pi_1(Y^\circ)\) therefore lands in \(Q\), giving the
commutative diagram
\[
\begin{array}{ccc}
 \pi_1(\widetilde S^\circ)&\longrightarrow&\pi_1(Y^\circ)\\
 \big\downarrow&&\big\downarrow\psi_Y\\
 \pi_1(S)&\xrightarrow{\nu}&N/B_0\ \subset\ G/B_0 .
\end{array}
\]
The left map is surjective, by removal of an analytic subset and
bimeromorphic invariance of the smooth fundamental group.
The map \(\nu\) is surjective by the definition of \(N\).
Thus \(\psi_Y\) is surjective. The upper map is the map of a
smooth proper fibration with connected fibres, so this is the
monodromy compatible with \(S\dashrightarrow Y\).

Transversality gives \(\Delta_Y=D_Z|_Y\) with simple normal
crossing support. For a deleted prime divisor \(D\), if \(E\) is a
component of \(D\cap Y\), its
meridian \(\mu_E\) maps to a conjugate of the meridian \(\mu_D\)
in \(Z^\circ\). The injection \(Q\subset R\) therefore gives
\[
 \operatorname{ord}_Q\bigl(\psi_Y(\mu_E)\bigr)
 =\operatorname{ord}_R\bigl(\psi_Z(\mu_D)\bigr)=m_D.
\]
The proper fibre-dimension theorem ensures that the deleted subsets
of codimension at least two remain of codimension at least two on
this very general restriction, and
filling them does not change the fundamental group. Filling
divisors of order one kills already trivial meridians. The
\(m_D\)-th powers of the remaining meridians are also killed,
so \(\psi_Y\) factors as a surjection
\[
 \pi_1^\orb(Y,\Delta_Y)\twoheadrightarrow Q.
\]

Finally fix a finite-index subgroup \(Q_0\le Q\), and let \(N_0\)
be its inverse image in \(N\). Lemma~\ref{lem:realize-finite-index}
gives a further global finite cover with group \(G'\) such that
\[
 N'=N\cap G'\subseteq N_0,\qquad
 Q'=\im(N'\to Q)\cong N'/(N'\cap B_0)\subseteq Q_0.
\]
The image \(Q'\) has finite index in \(Q\). Choose a point of the
new Albanese torus above the current \(a\). Its fibre \(S'\) is a
connected component of the pullback of \(S\), because the new
Albanese fibres are connected and the map of tori is finite.
Thus \(S'\to S\) is a connected finite \'etale cover, and its
induced homomorphism \(\nu'\) has image \(Q'\). The good-point
convention makes \(S'\) eligible for the same minimum.

Choose this cover before choosing a further quotient. For every
surjection \(\psi:Q'\twoheadrightarrow P\) with \(P\) infinite,
the \(\Gamma\)-reduction of \(\psi\circ\nu'\) on \(S'\) has
resolved general fibre image \(C'\triangleleft N'\) with finite
image in \(P\). If \(N'/C'\) were finite, the whole image in
\(P\) would be a finite union of translates of that finite image,
which is impossible. The reduction is therefore an admissible
fibration in the minimum, and its dimension is at least \(\ell\).
This one cover works for every infinite quotient of \(Q'\), as
required by \eqref{hi:minimality}. The prepared model was arbitrary;
repeating the choice of a very general good \(S\) on each model
proves the assertion on every prepared model.
\end{proof}

\subsection{The two representation cases}

Apply Proposition~\ref{prop:minimal-base}. The quotient \(Q\)
is FAb by Proposition~\ref{prop:fibre-fab}. Suppose first that
all finite-dimensional complex linear images of \(Q\) are
finite. On each prepared model of \(Z\), choose the very general
good fibre and the data in Proposition~\ref{prop:minimal-base}.
They satisfy the hypotheses of
Theorem~\ref{thm:hilbert-positivity}. Hence
\[
 K_Y+\Delta_Y\quad\text{is big}
\]
on the very general torus fibre, of positive dimension \(\ell\).
The same minimum applies after the fresh fibre choice on each
prepared model. Theorem~\ref{thm:torus-descent} contradicts the
specialness of \(X\).

Suppose instead that \(Q\) has an infinite complex linear image.
An infinite virtually solvable image is impossible for an FAb
group, as observed in Section~\ref{sec:preliminaries}. Pass to
finite index so that the Zariski closure \(\mathbf L\) is connected.
Its semisimple quotient is nontrivial. Choose a simple factor
\(\mathbf H\) of the adjoint semisimple quotient
\[
 \bigl(\mathbf L/\operatorname{Rad}(\mathbf L)\bigr)_{\mathrm{ad}},
\]
where \(\operatorname{Rad}(\mathbf L)\) is the solvable radical.
This quotient is a product of connected simple adjoint groups. Projection
gives a Zariski-dense image in \(\mathbf H\). In its faithful
adjoint representation this image is finitely generated and linear,
so Selberg's lemma \cite{Selberg1960} gives a finite-index subgroup
with torsion-free image. Pull this condition back to \(N\) and
realize it on a global finite cover by
Lemma~\ref{lem:realize-finite-index}. Write \(X',G',N',A'\) for
the resulting cover and its groups and Albanese torus. The induced
\(\rho_N:N'\to\mathbf H(\C)\) has torsion-free image. That image
is still Zariski dense, since it has finite index in a dense subgroup
of the connected group \(\mathbf H\).

Take the Stein factorization \(c':X'\dashrightarrow Z'\) of the
old map \(c\) after this cover. A connected \(c'\)-fibre maps under
the new Albanese map into a finite fibre of \(A'\to A\), and hence
to one point. Thus \(c'\) lies over a map \(g':Z'\to A'\);
connected Albanese fibres again make \(g'\) a fibration.
Use a smooth model for \(Z'\). Its relative dimension is still \(\ell\):
the new Stein base is finite over the old base on the general locus,
and \(A'\to A\) is an isogeny. Let \(B_0'\triangleleft G'\) be the
actual image of a general \(c'\)-fibre group. It lies in \(N'\).
A general \(c'\)-fibre is a connected component of a finite cover
of an old \(c\)-fibre, so its image in \(G\) lies in \(B_0\).
The representation \(\rho_N\), defined through the old \(Q=N/B_0\),
therefore kills \(B_0'\). Set
\[
 R'=G'/B_0',\qquad
 \overline\rho:N'/B_0'\longrightarrow\mathbf H(\C).
\]
Albanese exactness gives
\[
 1\longrightarrow N'/B_0'\longrightarrow R'
 \longrightarrow\pi_1(A')\longrightarrow1.
\]

On each prepared model of \(Z'\), use the actual inertial boundary
for \(R'\) and choose a fresh very general good Albanese fibre \(S'\)
and corresponding torus fibre \(Y\). Denote the restricted boundary
by \(\Delta_Y\), and let \(Y^\circ\) be the restricted smooth
fibration locus. The monodromy argument in
Proposition~\ref{prop:minimal-base}, applied to this restricted
fibration, gives
\[
 \psi'_Y:\pi_1(Y^\circ)\twoheadrightarrow N'/B_0'\subset R'.
\]
Compose it with \(\overline\rho\). A meridian has finite order in
\(R'\); its image under \(\overline\rho\circ\psi'_Y\) has order
dividing that actual inertial order. Since the representation image
is torsion free, every meridian is killed. Filling the deleted locus
therefore gives a representation
\(\rho_Y:\pi_1(Y)\to\mathbf H(\C)\) of the ordinary fundamental
group, with the same torsion-free Zariski-dense image.

This representation is generically large. Otherwise its
\(\Gamma\)-reduction has base dimension less than
\(\dim Y=\ell\). Compose it with \(S'\dashrightarrow Y\) and take
connected fibres by Stein factorization. A resolved general fibre
maps to a resolved \(\Gamma\)-fibre, so its group image
\(C'\triangleleft N'\) has finite image under \(\rho_N\).
The total image is infinite; if \(N'/C'\) were finite, the total
representation image would be a finite union of translates of that
finite fibre image.
Thus \(N'/C'\) is infinite, and this composed fibration would be
admissible in the minimum with base dimension less than \(\ell\),
a contradiction.

By Theorem~\ref{thm:linear-general-type}, \(K_Y\) is big, and hence
so is \(K_Y+\Delta_Y\). The argument applies after the fresh good
fibre choice on every prepared model, using the same minimum and
the boundary defined by the fixed actual quotient \(R'\) on each model.
Theorem~\ref{thm:torus-descent} yields the same contradiction.

Both cases exclude infinite \(N\). Thus \(N\) is finite and
the exact sequence of Lemma~\ref{lem:albanese-exactness}, followed
by Lemma~\ref{lem:finite-by-abelian}, makes \(G\) virtually
abelian. The group of the finite cover used at the start has
finite index in the original fundamental group, so the same
conclusion holds for the original \(X\). A finite-index
subgroup corresponds to a connected finite topological cover;
the complex structure lifts uniquely, making it an unramified
holomorphic cover. Base-point changes only conjugate the
subgroups. This completes the proof of
Theorem~\ref{thm:abelianity}.\qed
